%%%%%%%%%%%%%%%%%%%%%%%%%%%%%%%%%%%%%%%%%%%%%%%%%%%%%%%%%%%%%%%%%%%%%%%%%%%%%%%%
% CHAPTER 5
%%%%%%%%%%%%%%%%%%%%%%%%%%%%%%%%%%%%%%%%%%%%%%%%%%%%%%%%%%%%%%%%%%%%%%%%%%%%%%%%

\chapter{Continuous-Time Fourier Transform}

\begin{tcolorbox}[colback=blue!5!white,colframe=blue!70!black,title=Learning Objectives]

After studying this chapter, the reader will be able to

\begin{itemize}
\item Understand the need for the Fourier Transform.
\item Derive the Continuous-Time Fourier Transform (CTFT).
\item Compute CTFTs of common signals.
\item Understand the inverse Fourier Transform.
\item Interpret magnitude and phase spectra.
\item Use the duality between time and frequency domains.
\end{itemize}

\end{tcolorbox}

%%%%%%%%%%%%%%%%%%%%%%%%%%%%%%%%%%%%%%%%%%%%%%%%%%%%%%%%%%%%%%%%%%%%%%%%%%%%%%%%

\section{Introduction}

Fourier Series is applicable only to \textbf{periodic signals}. However, many
practical signals are \textbf{aperiodic}, such as

\begin{itemize}
\item speech segments,
\item radar pulses,
\item ECG pulses,
\item data packets,
\item and transient responses of circuits.
\end{itemize}

To analyze the frequency content of aperiodic signals, we use the
\textbf{Continuous-Time Fourier Transform (CTFT)}.

The Fourier Transform converts a signal from the

\[
\boxed{
\text{Time Domain }x(t)
\quad \longleftrightarrow \quad
\text{Frequency Domain }X(\omega)
}
\]

%%%%%%%%%%%%%%%%%%%%%%%%%%%%%%%%%%%%%%%%%%%%%%%%%%%%%%%%%%%%%%%%%%%%%%%%%%%%%%%%

\section{From Fourier Series to Fourier Transform}

Consider a periodic signal with period \(T\).

Its exponential Fourier series is

\[
x(t)=\sum_{n=-\infty}^{\infty}C_n e^{jn\omega_0 t}
\]

where

\[
\omega_0=\frac{2\pi}{T}.
\]

As \(T\to\infty\),

\begin{itemize}
\item the signal becomes aperiodic,
\item the spacing between harmonics approaches zero,
\item and the discrete spectrum becomes a continuous spectrum.
\end{itemize}

The summation changes into an integral, leading to the Fourier Transform.

%%%%%%%%%%%%%%%%%%%%%%%%%%%%%%%%%%%%%%%%%%%%%%%%%%%%%%%%%%%%%%%%%%%%%%%%%%%%%%%%

\section{Continuous-Time Fourier Transform}

%%%%%%%%%%%%%%%%%%%%%%%%%%%%%%%%%%%%%%%%%%%%%%%%%%%%%%%%%%%%%%%%%%%%%%%%%%%%%%%%

\subsection{Definition}

\[
\boxed{
X(\omega)=
\int_{-\infty}^{\infty}
x(t)e^{-j\omega t}\,dt
}
\]

where

\begin{itemize}
\item \(x(t)\): time-domain signal,
\item \(X(\omega)\): frequency-domain representation,
\item \(\omega\): angular frequency (rad/s).
\end{itemize}

%%%%%%%%%%%%%%%%%%%%%%%%%%%%%%%%%%%%%%%%%%%%%%%%%%%%%%%%%%%%%%%%%%%%%%%%%%%%%%%%

\section{Inverse Fourier Transform}

The original signal can be recovered using

\[
\boxed{
x(t)=
\frac{1}{2\pi}
\int_{-\infty}^{\infty}
X(\omega)e^{j\omega t}\,d\omega
}
\]

Together, these form the Fourier transform pair:

\[
\boxed
{
x(t)\ \xleftrightarrow{\mathcal{F}}\ X(\omega)
}
\]

%%%%%%%%%%%%%%%%%%%%%%%%%%%%%%%%%%%%%%%%%%%%%%%%%%%%%%%%%%%%%%%%%%%%%%%%%%%%%%%%

\section{Interpretation}

The exponential term

\[
e^{-j\omega t}
\]

acts as a basis function of frequency \(\omega\).

The transform coefficient \(X(\omega)\) tells us

\begin{itemize}
\item how much of frequency \(\omega\) is present,
\item and what phase that frequency component has.
\end{itemize}

%%%%%%%%%%%%%%%%%%%%%%%%%%%%%%%%%%%%%%%%%%%%%%%%%%%%%%%%%%%%%%%%%%%%%%%%%%%%%%%%

\section{Magnitude and Phase Spectrum}

Since \(X(\omega)\) is generally complex,

\[
X(\omega)=|X(\omega)|e^{j\angle X(\omega)}.
\]

%%%%%%%%%%%%%%%%%%%%%%%%%%%%%%%%%%%%%%%%%%%%%%%%%%%%%%%%%%%%%%%%%%%%%%%%%%%%%%%%

\subsection{Magnitude Spectrum}

\[
\boxed{
|X(\omega)|
}
\]

indicates the strength of each frequency component.

%%%%%%%%%%%%%%%%%%%%%%%%%%%%%%%%%%%%%%%%%%%%%%%%%%%%%%%%%%%%%%%%%%%%%%%%%%%%%%%%

\subsection{Phase Spectrum}

\[
\boxed{
\angle X(\omega)
}
\]

indicates the phase shift of each frequency component.

%%%%%%%%%%%%%%%%%%%%%%%%%%%%%%%%%%%%%%%%%%%%%%%%%%%%%%%%%%%%%%%%%%%%%%%%%%%%%%%%

\section{CTFT of the Unit Impulse}

Let

\[
x(t)=\delta(t).
\]

Using the definition,

\[
X(\omega)=
\int_{-\infty}^{\infty}
\delta(t)e^{-j\omega t}\,dt.
\]

Using the sifting property,

\[
\boxed{
X(\omega)=1
}
\]

Therefore,

\[
\boxed{
\delta(t)\ \xleftrightarrow{\mathcal{F}}\ 1
}
\]

This means the impulse contains all frequencies equally.

%%%%%%%%%%%%%%%%%%%%%%%%%%%%%%%%%%%%%%%%%%%%%%%%%%%%%%%%%%%%%%%%%%%%%%%%%%%%%%%%

\section{CTFT of a Constant}

Let

\[
x(t)=1.
\]

Then,

\[
X(\omega)=
\int_{-\infty}^{\infty}e^{-j\omega t}dt.
\]

This integral is not an ordinary function and is represented using the Dirac
delta function:

\[
\boxed{
1\ \xleftrightarrow{\mathcal{F}}\ 2\pi\delta(\omega)
}
\]

A constant signal contains only zero frequency (DC).

%%%%%%%%%%%%%%%%%%%%%%%%%%%%%%%%%%%%%%%%%%%%%%%%%%%%%%%%%%%%%%%%%%%%%%%%%%%%%%%%

\section{CTFT of an Exponential}

Consider

\[
x(t)=e^{-at}u(t),\qquad a>0.
\]

%%%%%%%%%%%%%%%%%%%%%%%%%%%%%%%%%%%%%%%%%%%%%%%%%%%%%%%%%%%%%%%%%%%%%%%%%%%%%%%%

\subsection{Computation}

\[
X(\omega)=
\int_0^\infty e^{-at}e^{-j\omega t}dt
\]

\[
=
\int_0^\infty e^{-(a+j\omega)t}dt
\]

\[
=
\left[
-\frac{e^{-(a+j\omega)t}}{a+j\omega}
\right]_0^\infty
\]

Since \(a>0\),

\[
e^{-(a+j\omega)t}\to0
\quad \text{as}\quad t\to\infty.
\]

Hence,

\[
\boxed{
X(\omega)=\frac{1}{a+j\omega}
}
\]

Therefore,

\[
\boxed{
e^{-at}u(t)
\ \xleftrightarrow{\mathcal{F}}\
\frac{1}{a+j\omega}
}
\]

%%%%%%%%%%%%%%%%%%%%%%%%%%%%%%%%%%%%%%%%%%%%%%%%%%%%%%%%%%%%%%%%%%%%%%%%%%%%%%%%

\section{Magnitude and Phase}

For

\[
X(\omega)=\frac{1}{a+j\omega},
\]

%%%%%%%%%%%%%%%%%%%%%%%%%%%%%%%%%%%%%%%%%%%%%%%%%%%%%%%%%%%%%%%%%%%%%%%%%%%%%%%%

\subsection{Magnitude}

\[
\boxed{
|X(\omega)|=
\frac{1}{\sqrt{a^2+\omega^2}}
}
\]

%%%%%%%%%%%%%%%%%%%%%%%%%%%%%%%%%%%%%%%%%%%%%%%%%%%%%%%%%%%%%%%%%%%%%%%%%%%%%%%%

\subsection{Phase}

\[
\boxed{
\angle X(\omega)=
-\tan^{-1}\left(\frac{\omega}{a}\right)
}
\]

%%%%%%%%%%%%%%%%%%%%%%%%%%%%%%%%%%%%%%%%%%%%%%%%%%%%%%%%%%%%%%%%%%%%%%%%%%%%%%%%

\section{Rectangular Pulse}

Consider

\[
x(t)=
\begin{cases}
1, & |t|<T/2\\
0, & |t|>T/2
\end{cases}
\]

%%%%%%%%%%%%%%%%%%%%%%%%%%%%%%%%%%%%%%%%%%%%%%%%%%%%%%%%%%%%%%%%%%%%%%%%%%%%%%%%

\subsection{Transform}

\[
X(\omega)=
\int_{-T/2}^{T/2}e^{-j\omega t}dt
\]

\[
=
\left[
\frac{e^{-j\omega t}}{-j\omega}
\right]_{-T/2}^{T/2}
\]

\[
=
\frac{2\sin(\omega T/2)}{\omega}
\]

Define the sinc function as

\[
\boxed{
\mathrm{sinc}(x)=\frac{\sin x}{x}
}
\]

Then,

\[
\boxed{
X(\omega)=
T\,\mathrm{sinc}\left(\frac{\omega T}{2}\right)
}
\]

%%%%%%%%%%%%%%%%%%%%%%%%%%%%%%%%%%%%%%%%%%%%%%%%%%%%%%%%%%%%%%%%%%%%%%%%%%%%%%%%

\section{Sinc Function}

The sinc function has

\begin{itemize}
\item a main lobe,
\item and infinitely many side lobes.
\end{itemize}

Zeros occur at

\[
\boxed{
\omega=\frac{2\pi n}{T},
\qquad n=\pm1,\pm2,\ldots
}
\]

%%%%%%%%%%%%%%%%%%%%%%%%%%%%%%%%%%%%%%%%%%%%%%%%%%%%%%%%%%%%%%%%%%%%%%%%%%%%%%%%

\section{Time-Domain Width vs Frequency-Domain Width}

A very important observation is

\[
\boxed{
\text{Narrow pulse in time}
\Rightarrow
\text{Wide spectrum in frequency}
}
\]

and

\[
\boxed{
\text{Wide pulse in time}
\Rightarrow
\text{Narrow spectrum in frequency}
}
\]

This is the \textbf{time-frequency tradeoff}.

%%%%%%%%%%%%%%%%%%%%%%%%%%%%%%%%%%%%%%%%%%%%%%%%%%%%%%%%%%%%%%%%%%%%%%%%%%%%%%%%

\section{Fourier Transform Pairs}

\begin{table}[H]
\centering
\caption{Important CTFT pairs}
\begin{tabular}{|p{5cm}|p{6cm}|}
\hline
\textbf{Time Domain} & \textbf{Frequency Domain} \\
\hline
\(\delta(t)\) & \(1\) \\
\hline
\(1\) & \(2\pi\delta(\omega)\) \\
\hline
\(e^{-at}u(t)\) & \(\dfrac{1}{a+j\omega}\) \\
\hline
Rectangular pulse & \(T\,\mathrm{sinc}(\omega T/2)\) \\
\hline
\end{tabular}
\end{table}

%%%%%%%%%%%%%%%%%%%%%%%%%%%%%%%%%%%%%%%%%%%%%%%%%%%%%%%%%%%%%%%%%%%%%%%%%%%%%%%%

\section{Existence of the Fourier Transform}

A sufficient condition is

\[
\boxed{
\int_{-\infty}^{\infty}|x(t)|dt<\infty
}
\]

Signals satisfying this are called \textbf{absolutely integrable}.

Examples:

\begin{itemize}
\item \(e^{-at}u(t)\) \(\rightarrow\) transform exists.
\item \(u(t)\) \(\rightarrow\) not absolutely integrable.
\end{itemize}

%%%%%%%%%%%%%%%%%%%%%%%%%%%%%%%%%%%%%%%%%%%%%%%%%%%%%%%%%%%%%%%%%%%%%%%%%%%%%%%%

\section{Physical Meaning}

The Fourier Transform tells us

\begin{itemize}
\item which frequencies are present,
\item how strong they are,
\item and how they are phase shifted.
\end{itemize}

For communication systems, the spectrum determines

\begin{itemize}
\item bandwidth,
\item filter requirements,
\item channel allocation,
\item and noise performance.
\end{itemize}

%%%%%%%%%%%%%%%%%%%%%%%%%%%%%%%%%%%%%%%%%%%%%%%%%%%%%%%%%%%%%%%%%%%%%%%%%%%%%%%%

\begin{tcolorbox}[title=Quick Recall,colback=yellow!10]

\[
X(\omega)=
\int_{-\infty}^{\infty}
x(t)e^{-j\omega t}dt
\]

\[
x(t)=
\frac{1}{2\pi}
\int_{-\infty}^{\infty}
X(\omega)e^{j\omega t}d\omega
\]

\[
\delta(t)\leftrightarrow 1
\]

\[
1\leftrightarrow 2\pi\delta(\omega)
\]

\[
e^{-at}u(t)\leftrightarrow \frac{1}{a+j\omega}
\]

\[
\text{Rectangular pulse}\leftrightarrow T\,\mathrm{sinc}\left(\frac{\omega T}{2}\right)
\]

\end{tcolorbox}
%%%%%%%%%%%%%%%%%%%%%%%%%%%%%%%%%%%%%%%%%%%%%%%%%%%%%%%%%%%%%%%%%%%%%%%%%%%%%%%%
%%%%%%%%%%%%%%%%%%%%%%%%%%%%%%%%%%%%%%%%%%%%%%%%%%%%%%%%%%%%%%%%%%%%%%%%%%%%%%%%
\section{Linearity Property}

If

\[
x_1(t)\xleftrightarrow{\mathcal{F}}X_1(\omega)
\]

and

\[
x_2(t)\xleftrightarrow{\mathcal{F}}X_2(\omega),
\]

then

\[
\boxed{
a_1x_1(t)+a_2x_2(t)
\xleftrightarrow{\mathcal{F}}
a_1X_1(\omega)+a_2X_2(\omega)
}
\]

This follows directly from the linearity of integration.

%%%%%%%%%%%%%%%%%%%%%%%%%%%%%%%%%%%%%%%%%%%%%%%%%%%%%%%%%%%%%%%%%%%%%%%%%%%%%%%%

\section{Time-Shifting Property}

Let

\[
x(t)\xleftrightarrow{\mathcal{F}}X(\omega).
\]

Then a delay of \(t_0\) seconds gives

\[
\boxed{
x(t-t_0)
\xleftrightarrow{\mathcal{F}}
e^{-j\omega t_0}X(\omega)
}
\]

%%%%%%%%%%%%%%%%%%%%%%%%%%%%%%%%%%%%%%%%%%%%%%%%%%%%%%%%%%%%%%%%%%%%%%%%%%%%%%%%

\subsection{Interpretation}

\begin{itemize}
\item Time delay does not change the magnitude spectrum.
\item It introduces a linear phase shift.
\end{itemize}

%%%%%%%%%%%%%%%%%%%%%%%%%%%%%%%%%%%%%%%%%%%%%%%%%%%%%%%%%%%%%%%%%%%%%%%%%%%%%%%%

\subsection{Example}

\[
e^{-a(t-2)}u(t-2)
\xleftrightarrow{\mathcal{F}}
e^{-j2\omega}\frac{1}{a+j\omega}
\]

%%%%%%%%%%%%%%%%%%%%%%%%%%%%%%%%%%%%%%%%%%%%%%%%%%%%%%%%%%%%%%%%%%%%%%%%%%%%%%%%

\section{Frequency-Shifting Property}

If

\[
x(t)\xleftrightarrow{\mathcal{F}}X(\omega),
\]

then

\[
\boxed{
x(t)e^{j\omega_0 t}
\xleftrightarrow{\mathcal{F}}
X(\omega-\omega_0)
}
\]

Multiplication by a complex exponential shifts the spectrum.

%%%%%%%%%%%%%%%%%%%%%%%%%%%%%%%%%%%%%%%%%%%%%%%%%%%%%%%%%%%%%%%%%%%%%%%%%%%%%%%%

\subsection{Cosine Modulation}

Using

\[
\cos\omega_0 t=
\frac{1}{2}
\left(
e^{j\omega_0 t}+e^{-j\omega_0 t}
\right),
\]

we get

\[
\boxed{
x(t)\cos\omega_0 t
\xleftrightarrow{\mathcal{F}}
\frac{1}{2}
\left[
X(\omega-\omega_0)+X(\omega+\omega_0)
\right]
}
\]

This is the basis of amplitude modulation.

%%%%%%%%%%%%%%%%%%%%%%%%%%%%%%%%%%%%%%%%%%%%%%%%%%%%%%%%%%%%%%%%%%%%%%%%%%%%%%%%

\section{Time-Scaling Property}

For \(a\neq0\),

\[
\boxed{
x(at)
\xleftrightarrow{\mathcal{F}}
\frac{1}{|a|}X\left(\frac{\omega}{a}\right)
}
\]

%%%%%%%%%%%%%%%%%%%%%%%%%%%%%%%%%%%%%%%%%%%%%%%%%%%%%%%%%%%%%%%%%%%%%%%%%%%%%%%%

\subsection{Interpretation}

\begin{itemize}
\item Compressing a signal in time expands its spectrum.
\item Expanding a signal in time compresses its spectrum.
\end{itemize}

%%%%%%%%%%%%%%%%%%%%%%%%%%%%%%%%%%%%%%%%%%%%%%%%%%%%%%%%%%%%%%%%%%%%%%%%%%%%%%%%

\subsection{Example}

\[
x(2t)
\xleftrightarrow{\mathcal{F}}
\frac{1}{2}X\left(\frac{\omega}{2}\right)
\]

%%%%%%%%%%%%%%%%%%%%%%%%%%%%%%%%%%%%%%%%%%%%%%%%%%%%%%%%%%%%%%%%%%%%%%%%%%%%%%%%

\section{Differentiation in Time}

If

\[
x(t)\xleftrightarrow{\mathcal{F}}X(\omega),
\]

then

\[
\boxed{
\frac{dx(t)}{dt}
\xleftrightarrow{\mathcal{F}}
j\omega X(\omega)
}
\]

%%%%%%%%%%%%%%%%%%%%%%%%%%%%%%%%%%%%%%%%%%%%%%%%%%%%%%%%%%%%%%%%%%%%%%%%%%%%%%%%

\subsection{Higher Derivatives}

\[
\boxed{
\frac{d^n x(t)}{dt^n}
\xleftrightarrow{\mathcal{F}}
(j\omega)^nX(\omega)
}
\]

Differentiation emphasizes high-frequency components.

%%%%%%%%%%%%%%%%%%%%%%%%%%%%%%%%%%%%%%%%%%%%%%%%%%%%%%%%%%%%%%%%%%%%%%%%%%%%%%%%

\section{Integration in Time}

\[
\boxed{
\int_{-\infty}^{t}x(\tau)\,d\tau
\xleftrightarrow{\mathcal{F}}
\frac{X(\omega)}{j\omega}+\pi X(0)\delta(\omega)
}
\]

Integration acts as a low-pass operation.

%%%%%%%%%%%%%%%%%%%%%%%%%%%%%%%%%%%%%%%%%%%%%%%%%%%%%%%%%%%%%%%%%%%%%%%%%%%%%%%%

\section{Convolution Property}

One of the most important properties is

\[
\boxed{
x(t)*h(t)
\xleftrightarrow{\mathcal{F}}
X(\omega)H(\omega)
}
\]

Thus, convolution in time becomes multiplication in frequency.

%%%%%%%%%%%%%%%%%%%%%%%%%%%%%%%%%%%%%%%%%%%%%%%%%%%%%%%%%%%%%%%%%%%%%%%%%%%%%%%%

\subsection{LTI System Representation}

For an LTI system,

\[
y(t)=x(t)*h(t).
\]

Taking Fourier transforms,

\[
\boxed{
Y(\omega)=X(\omega)H(\omega)
}
\]

where \(H(\omega)\) is the frequency response of the system.

%%%%%%%%%%%%%%%%%%%%%%%%%%%%%%%%%%%%%%%%%%%%%%%%%%%%%%%%%%%%%%%%%%%%%%%%%%%%%%%%

\section{Multiplication Property}

Multiplication in time corresponds to convolution in frequency:

\[
\boxed{
x(t)h(t)
\xleftrightarrow{\mathcal{F}}
\frac{1}{2\pi}X(\omega)*H(\omega)
}
\]

%%%%%%%%%%%%%%%%%%%%%%%%%%%%%%%%%%%%%%%%%%%%%%%%%%%%%%%%%%%%%%%%%%%%%%%%%%%%%%%%

\section{Duality Property}

If

\[
x(t)\xleftrightarrow{\mathcal{F}}X(\omega),
\]

then

\[
\boxed{
X(t)\xleftrightarrow{\mathcal{F}}2\pi x(-\omega)
}
\]

Duality allows new transform pairs to be obtained from known ones.

%%%%%%%%%%%%%%%%%%%%%%%%%%%%%%%%%%%%%%%%%%%%%%%%%%%%%%%%%%%%%%%%%%%%%%%%%%%%%%%%

\section{Parseval’s Theorem}

For energy signals,

\[
E=
\int_{-\infty}^{\infty}|x(t)|^2dt.
\]

In the frequency domain,

\[
\boxed{
E=
\frac{1}{2\pi}
\int_{-\infty}^{\infty}|X(\omega)|^2d\omega
}
\]

This shows that signal energy is conserved across domains.

%%%%%%%%%%%%%%%%%%%%%%%%%%%%%%%%%%%%%%%%%%%%%%%%%%%%%%%%%%%%%%%%%%%%%%%%%%%%%%%%

\section{Energy Spectral Density}

Define

\[
\boxed{
\Psi(\omega)=|X(\omega)|^2
}
\]

called the energy spectral density (ESD).

Then

\[
E=
\frac{1}{2\pi}
\int_{-\infty}^{\infty}\Psi(\omega)d\omega.
\]

%%%%%%%%%%%%%%%%%%%%%%%%%%%%%%%%%%%%%%%%%%%%%%%%%%%%%%%%%%%%%%%%%%%%%%%%%%%%%%%%

\section{Solved Example 5.1}

Find the CTFT of

\[
x(t)=e^{-2t}u(t).
\]

%%%%%%%%%%%%%%%%%%%%%%%%%%%%%%%%%%%%%%%%%%%%%%%%%%%%%%%%%%%%%%%%%%%%%%%%%%%%%%%%

\subsection{Solution}

Using the standard pair,

\[
\boxed{
X(\omega)=\frac{1}{2+j\omega}
}
\]

Magnitude:

\[
|X(\omega)|=
\frac{1}{\sqrt{4+\omega^2}}
\]

Phase:

\[
\angle X(\omega)=
-\tan^{-1}\left(\frac{\omega}{2}\right)
\]

%%%%%%%%%%%%%%%%%%%%%%%%%%%%%%%%%%%%%%%%%%%%%%%%%%%%%%%%%%%%%%%%%%%%%%%%%%%%%%%%

\section{Solved Example 5.2}

Find the CTFT of

\[
x(t)=u(t-3).
\]

%%%%%%%%%%%%%%%%%%%%%%%%%%%%%%%%%%%%%%%%%%%%%%%%%%%%%%%%%%%%%%%%%%%%%%%%%%%%%%%%

\subsection{Step 1}

We know

\[
u(t)\xleftrightarrow{\mathcal{F}}
\pi\delta(\omega)+\frac{1}{j\omega}.
\]

%%%%%%%%%%%%%%%%%%%%%%%%%%%%%%%%%%%%%%%%%%%%%%%%%%%%%%%%%%%%%%%%%%%%%%%%%%%%%%%%

\subsection{Step 2}

Apply time shift:

\[
\boxed{
u(t-3)
\xleftrightarrow{\mathcal{F}}
e^{-j3\omega}
\left(
\pi\delta(\omega)+\frac{1}{j\omega}
\right)
}
\]

%%%%%%%%%%%%%%%%%%%%%%%%%%%%%%%%%%%%%%%%%%%%%%%%%%%%%%%%%%%%%%%%%%%%%%%%%%%%%%%%

\section{Solved Example 5.3}

Find the CTFT of

\[
x(t)=\cos\omega_0 t.
\]

%%%%%%%%%%%%%%%%%%%%%%%%%%%%%%%%%%%%%%%%%%%%%%%%%%%%%%%%%%%%%%%%%%%%%%%%%%%%%%%%

\subsection{Solution}

Using

\[
\cos\omega_0 t=
\frac{1}{2}
\left(
e^{j\omega_0 t}+e^{-j\omega_0 t}
\right),
\]

and

\[
e^{j\omega_0 t}
\xleftrightarrow{\mathcal{F}}
2\pi\delta(\omega-\omega_0),
\]

we obtain

\[
\boxed{
\cos\omega_0 t
\xleftrightarrow{\mathcal{F}}
\pi
\left[
\delta(\omega-\omega_0)+\delta(\omega+\omega_0)
\right]
}
\]

A cosine has two spectral lines at \(\pm\omega_0\).

%%%%%%%%%%%%%%%%%%%%%%%%%%%%%%%%%%%%%%%%%%%%%%%%%%%%%%%%%%%%%%%%%%%%%%%%%%%%%%%%

\section{Solved Example 5.4}

Find the CTFT of a rectangular pulse

\[
x(t)=
\begin{cases}
1, & |t|<1\\
0, & |t|>1
\end{cases}
\]

%%%%%%%%%%%%%%%%%%%%%%%%%%%%%%%%%%%%%%%%%%%%%%%%%%%%%%%%%%%%%%%%%%%%%%%%%%%%%%%%

\subsection{Solution}

Here \(T=2\).

\[
X(\omega)=
\int_{-1}^{1}e^{-j\omega t}dt
\]

\[
=
\frac{2\sin\omega}{\omega}
\]

Hence,

\[
\boxed{
X(\omega)=2\,\mathrm{sinc}(\omega)
}
\]

with \(\mathrm{sinc}(\omega)=\dfrac{\sin\omega}{\omega}\).

%%%%%%%%%%%%%%%%%%%%%%%%%%%%%%%%%%%%%%%%%%%%%%%%%%%%%%%%%%%%%%%%%%%%%%%%%%%%%%%%

\section{Bandwidth}

The bandwidth of a signal is the range of frequencies over which its spectrum is significant.

%%%%%%%%%%%%%%%%%%%%%%%%%%%%%%%%%%%%%%%%%%%%%%%%%%%%%%%%%%%%%%%%%%%%%%%%%%%%%%%%

\subsection{Rectangular Pulse}

For a pulse of width \(T\),

\[
X(\omega)=
T\,\mathrm{sinc}\left(\frac{\omega T}{2}\right).
\]

The first zeros occur at

\[
\omega=\pm\frac{2\pi}{T}.
\]

Hence, the main-lobe bandwidth is approximately

\[
\boxed{
B\approx\frac{4\pi}{T}
}
\]

Narrower pulses require larger bandwidth.

%%%%%%%%%%%%%%%%%%%%%%%%%%%%%%%%%%%%%%%%%%%%%%%%%%%%%%%%%%%%%%%%%%%%%%%%%%%%%%%%

\section{Summary of CTFT Properties}

\begin{table}[H]
\centering
\caption{Important CTFT properties}
\begin{tabular}{|p{5cm}|p{6cm}|}
\hline
\textbf{Time Domain} & \textbf{Frequency Domain} \\
\hline
\(a_1x_1+a_2x_2\) & \(a_1X_1+a_2X_2\) \\
\hline
\(x(t-t_0)\) & \(e^{-j\omega t_0}X(\omega)\) \\
\hline
\(x(t)e^{j\omega_0 t}\) & \(X(\omega-\omega_0)\) \\
\hline
\(x(at)\) & \(\dfrac{1}{|a|}X(\omega/a)\) \\
\hline
\(\dfrac{dx}{dt}\) & \(j\omega X(\omega)\) \\
\hline
\(\int x(\tau)d\tau\) & \(\dfrac{X(\omega)}{j\omega}+\pi X(0)\delta(\omega)\) \\
\hline
\(x*h\) & \(XH\) \\
\hline
\(xh\) & \(\dfrac{1}{2\pi}X*H\) \\
\hline
\end{tabular}
\end{table}

%%%%%%%%%%%%%%%%%%%%%%%%%%%%%%%%%%%%%%%%%%%%%%%%%%%%%%%%%%%%%%%%%%%%%%%%%%%%%%%%

\section{Problem-Solving Strategy}

\subsection{Step 1}

Check whether the signal matches a standard transform pair.

%%%%%%%%%%%%%%%%%%%%%%%%%%%%%%%%%%%%%%%%%%%%%%%%%%%%%%%%%%%%%%%%%%%%%%%%%%%%%%%%

\subsection{Step 2}

Apply properties such as time shift, scaling or modulation.

%%%%%%%%%%%%%%%%%%%%%%%%%%%%%%%%%%%%%%%%%%%%%%%%%%%%%%%%%%%%%%%%%%%%%%%%%%%%%%%%

\subsection{Step 3}

Simplify the expression algebraically.

%%%%%%%%%%%%%%%%%%%%%%%%%%%%%%%%%%%%%%%%%%%%%%%%%%%%%%%%%%%%%%%%%%%%%%%%%%%%%%%%

\subsection{Step 4}

Find magnitude and phase if required.

%%%%%%%%%%%%%%%%%%%%%%%%%%%%%%%%%%%%%%%%%%%%%%%%%%%%%%%%%%%%%%%%%%%%%%%%%%%%%%%%

\subsection{Step 5}

For LTI systems, use

\[
Y(\omega)=X(\omega)H(\omega).
\]

%%%%%%%%%%%%%%%%%%%%%%%%%%%%%%%%%%%%%%%%%%%%%%%%%%%%%%%%%%%%%%%%%%%%%%%%%%%%%%%%

\section{Common GATE Mistakes}

\begin{enumerate}
\item Forgetting the factor \(\dfrac{1}{2\pi}\) in the inverse transform.
\item Confusing frequency shift with time shift.
\item Writing \(X(\omega-\omega_0)\) instead of \(X(\omega+\omega_0)\).
\item Forgetting the scaling factor \(\dfrac{1}{|a|}\).
\item Using ordinary multiplication instead of convolution.
\item Forgetting the delta functions in the transform of sinusoids.
\item Using \(\sin x/x\) without defining the sinc function.
\end{enumerate}

%%%%%%%%%%%%%%%%%%%%%%%%%%%%%%%%%%%%%%%%%%%%%%%%%%%%%%%%%%%%%%%%%%%%%%%%%%%%%%%%

\section{Chapter Summary}

In this chapter, we developed the Continuous-Time Fourier Transform for aperiodic signals and studied its inverse transform, magnitude spectrum and phase spectrum.

We derived the transforms of

\begin{itemize}
\item the impulse,
\item constant signal,
\item exponential signal,
\item and rectangular pulse.
\end{itemize}

We also studied the most important CTFT properties, including time shifting, frequency shifting, scaling, differentiation, integration, convolution and duality.

Finally, we introduced Parseval’s theorem and energy spectral density, which connect the time-domain energy of a signal to its frequency-domain representation.

%%%%%%%%%%%%%%%%%%%%%%%%%%%%%%%%%%%%%%%%%%%%%%%%%%%%%%%%%%%%%%%%%%%%%%%%%%%%%%%%

\begin{tcolorbox}[title=One-Minute Revision,colback=blue!5]

\[
X(\omega)=
\int_{-\infty}^{\infty}x(t)e^{-j\omega t}dt
\]

\[
x(t)=
\frac{1}{2\pi}
\int_{-\infty}^{\infty}X(\omega)e^{j\omega t}d\omega
\]

\[
x(t-t_0)\leftrightarrow e^{-j\omega t_0}X(\omega)
\]

\[
x(t)e^{j\omega_0 t}\leftrightarrow X(\omega-\omega_0)
\]

\[
x(at)\leftrightarrow \frac{1}{|a|}X\left(\frac{\omega}{a}\right)
\]

\[
\frac{dx}{dt}\leftrightarrow j\omega X(\omega)
\]

\[
x*h\leftrightarrow XH
\]

\[
E=
\frac{1}{2\pi}
\int_{-\infty}^{\infty}|X(\omega)|^2d\omega
\]

\[
e^{-at}u(t)\leftrightarrow \frac{1}{a+j\omega}
\]

\[
\text{Rectangular pulse}\leftrightarrow
T\,\mathrm{sinc}\left(\frac{\omega T}{2}\right)
\]

\end{tcolorbox}
%%%%%%%%%%%%%%%%%%%%%%%%%%%%%%%%%%%%%%%%%%%%%%%%%%%%%%%%%%%%%%%%%%%%%%%%%%%%%%%%